
\documentclass[aspectratio=169]{beamer}

\mode<presentation>
\usetheme{Boadilla}
\definecolor{NTHU1}{RGB}{127,16,132}
\definecolor{NTHU2}{RGB}{228,210,231}
\definecolor{blue}{RGB}{30,90,205}
\definecolor{red}{RGB}{213,94,0}
\definecolor{green}{RGB}{0,128,0}
\definecolor{charcoalgray}{RGB}{108,104,104}
\setbeamercolor{title}{fg=NTHU1}
\setbeamercolor{frametitle}{fg=NTHU1}
\setbeamercolor{block title}{bg=NTHU1, fg=white}
\setbeamercolor{block body}{bg=NTHU2}
\setbeamercolor{structure}{fg=NTHU1}
\setbeamercolor{item projected}{fg=white}
\setbeamercolor{item}{fg=NTHU1}
\setbeamercolor{subitem}{fg=NTHU1}
\setbeamercolor{section in toc}{fg=NTHU1}
\setbeamercolor{description item}{fg=NTHU1}
\setbeamercolor{caption name}{fg=NTHU1}
\setbeamercolor{button}{bg=NTHU1, fg=white}
\setbeamercolor{caption name}{fg=NTHU1}
\usepackage{graphics}
\usepackage{tikz}
\usepackage{amsmath}
\usepackage{bbm}
\usetikzlibrary{decorations.pathreplacing}
\usepackage{geometry}
\usepackage{booktabs}
\usepackage{bookmark}
\usepackage{multirow, makecell}
\usepackage{float}
\usepackage{fancyvrb}
\usepackage{caption}
\captionsetup[figure]{singlelinecheck = false, format= hang, justification=centering}
\captionsetup[subfigure]{singlelinecheck = false, format= hang, justification=raggedright}
\usepackage{subcaption}
\usepackage{adjustbox}
\usepackage{hyperref}
\usepackage{threeparttable}
\usepackage{appendixnumberbeamer} 
\usepackage[T1]{fontenc}
\usepackage[default]{lato} 
\usepackage{smartdiagram}
\smartdiagramset{
  back arrow disabled = true,
  module minimum width=1cm,
  module x sep = 3,
  uniform arrow color=true,
}
\usepackage{xcolor}
\usepackage{colortbl}
	\definecolor{light-gray}{gray}{0.99}

\newenvironment{wideitemize}{\itemize\addtolength{\itemsep}{10pt}}{\enditemize}
\newenvironment{wideenumerate}{\enumerate\addtolength{\itemsep}{10pt}}{\endenumerate}
\newenvironment{widedescription}{\description\addtolength{\itemsep}{10pt}}{\enddescription}
\hypersetup{
colorlinks=true,
linkcolor=NTHU1,
filecolor=green, 
urlcolor=blue,
}
\beamertemplatenavigationsymbolsempty
\setbeamercolor{author in head/foot}{bg=white, fg=NTHU1}
\setbeamercolor{title in head/foot}{bg=white, fg=NTHU1}
\setbeamercolor{date in head/foot}{bg=white, fg=NTHU1}
\setbeamercolor{section in head/foot}{bg=white, fg=NTHU1}
\setbeamercolor{headline}{bg=white}
\setbeamertemplate{footline}{
    \leavevmode%
    \hbox{%
        \begin{beamercolorbox}[wd=.333333\paperwidth,ht=2.25ex,dp=1ex,center]{date in head/foot}%
            \usebeamerfont{date in head/foot}%\insertdate
        \end{beamercolorbox}%
        \begin{beamercolorbox}[wd=.333333\paperwidth,ht=2.25ex,dp=1ex,center]{title in head/foot}%
            \usebeamerfont{title in head/foot}\insertshorttitle
        \end{beamercolorbox}%
        \begin{beamercolorbox}[wd=.333333\paperwidth,ht=2.25ex,dp=1ex,right]{date in head/foot}%
            \usebeamerfont{date in head/foot}\insertframenumber{}\hspace*{2em}
        \end{beamercolorbox}}%
        \vskip0pt%
    }
    %% May need to script the introduction!!!!!!
%\setbeamercolor{page number in head/foot}{fg=NTHU1}
\setbeamertemplate{section in toc}[sections numbered]
\setbeamertemplate{subsection in toc}{\leavevmode\leftskip=3em\rlap{\hskip-1.75em\inserttocsectionnumber.\inserttocsubsectionnumber}
\inserttocsubsection\par}
\setbeamerfont{subsection in toc}{size=\footnotesize}
%\setbeamertemplate{headline}{%
  %\begin{beamercolorbox}[ht=5.5ex]{section in head/foot}
    %\vskip2pt\insertnavigation{0.33\paperwidth}\vskip2pt
  %\end{beamercolorbox}%
%}
\newenvironment{transitionframe}{\setbeamercolor{background canvas}{bg=white}\setbeamertemplate{footline}{} \begin{frame}[noframenumbering]}{\end{frame}}


\makeatletter
\let\@@magyar@captionfix\relax
\makeatother


\title[]{Public Economics and Development} % Change this regularly
\subtitle[]{Week 9: Targeting beneficiaries of public goods and services}
\author[Lee]{Seung-hun Lee }
\institute[]{Taipei School of Economics\\ National Tsing Hua University}

\date[]{April 21st, 2026}

\begin{document}
\begin{transitionframe}
\titlepage
\end{transitionframe}


%%% Color slides for section headers: Use for colloquium version (The ones bewteen \iffals and \fi)

\section{Overview setup of the topic}

\begin{frame}
\frametitle{Social Protection: Scale and Composition}

Growth of social protection programs (ILO World Social Protection Report 2024-26)

\begin{itemize}\setlength\itemsep{4pt}
    \item \textbf{Old-age pensions}: 79.6\% of people above retirement age receive a pension
    \item \textbf{Disability benefits}: 38.9\% of people with severe disabilities receive a benefit
    \item \textbf{Child/family benefits}: 28.2\% of children globally receive a family or child benefit
    \item \textbf{Unemployment benefits}: only 16.7\% of unemployed people receive cash benefits 
    \item \textbf{Others:} Insurance programs against climate change, Income support etc..
\end{itemize}\medskip

Overall, \textbf{52.4\%} of the world's population are covered by at least one program
\begin{itemize}\setlength\itemsep{4pt}  
\item But \textbf{3.8 billion people} remain uncovered
  \item Most of the uncovered are in low-income countries
  \begin{itemize}\setlength\itemsep{4pt}
  \item Social protection spending averages only 3.7\% of GDP vs.\ 12.9\% globally.
\end{itemize}
 \end{itemize}
\end{frame}

\begin{frame}
\frametitle{Variation in coverage under cash benefit programs}
\begin{figure}
  \centering
  \includegraphics[keepaspectratio, width=0.7\textwidth]{fig1.png}
\end{figure}
\end{frame}


\begin{frame}
\frametitle{Variation in population covered by social progreams}
\begin{figure}
  \centering
  \includegraphics[keepaspectratio, width=1\textwidth]{fig2.png}
\end{figure}
\end{frame}


\begin{frame}
\frametitle{The Targeting vs.\ Universalism Debate}

The case for targeting:
\begin{itemize}\setlength\itemsep{3pt}
    \item Cost-effective: Concentrates resources on those who need them most 
    \item Fiscal constraints make universal programs hard to sustain
\end{itemize}\medskip

The case for universalism:
\begin{itemize}\setlength\itemsep{3pt}
    \item Identifying the poor can be administratively costly and error-prone
    \item Standard measures of poverty can be uninformative about actual welfare
\end{itemize}\medskip

Reality: Theoretical advantage of targeting can erode in practice
\begin{itemize}\setlength\itemsep{3pt}
    \item Corruption and misallocation: Intended beneficiaries may not receive benefits
    \item Low-take up due to information gaps, stigma, or administrative barriers
\end{itemize}

\end{frame}



\begin{frame}
\frametitle{Why Don't Eligible Households Take Up Benefits? }

Why take up is not complete among eligibles (Currie 2006)\medskip

\begin{itemize}\setlength\itemsep{6pt}
    \item \textbf{Stigma}: ``Shame'' of receiving benefits, 
    \item \textbf{Information barriers}: Those legally eligible unaware of their status or the procedure
    \item \textbf{Transaction costs}: Paperworks, or going to offices, can be costly and time-consuming 
    \item \textbf{Social networks}: Fewer connections to prior recipients $\to$ less likely to participate
\end{itemize}\medskip

These costs may fall disproportionately on the neediest households
\begin{itemize}\setlength\itemsep{6pt}
\item Maybe even more time and cost constrained than typical households
\item Could face higher barriers to the same information
\end{itemize} \medskip

$\Rightarrow$ Even a well-designed targeting rule may fail to reach its intended beneficiaries if these costs are ignored. \textit{So how can we lower them?}

\end{frame}

\begin{transitionframe}
\frametitle{Roadmap for today}
\tableofcontents
\end{transitionframe}








\section{Theories on targeting transfers}
\subsection{Nichols, Zeckhauser (1982): Targeting Transfers through Restrictions on Recipients}
\begin{transitionframe}
\begin{center}
  \begin{huge}
    \textcolor{NTHU1}{%
Nichols, Zeckhauser (1982):\\ Targeting Transfers through Restrictions on Recipients
    }
  \end{huge}
\end{center}
\end{transitionframe}


\begin{frame}
\frametitle{Q: How to model social programs in presence of unobservable traits? } 
Government wants to redistribute to maximize social welfare
\begin{itemize}\setlength\itemsep{3pt}
\item Ideally, we base lump-sum tax based on random (and un-distortable) characteristics
\item In practice, only the partial indicators of such traits are observed
\begin{itemize}\setlength\itemsep{3pt}
\item Earnings: Is it because individuals are productive or they simply work long hours?
\item Expenditure: Reflection of true preference or other factors leading up to purchase?
\end{itemize}
\item In such situations, we resort to second-best solutions using observable instruments
\end{itemize}\medskip
This paper provides a model on what instruments to use to target redistribution programs
\begin{itemize}\setlength\itemsep{3pt}
\item Target efficiency vs productive efficiency
\item Introduces model for restrictrions: Means-tested programs, in-kind transfers etc
\item Compare its performance to programs solely relying on income tax and cash transfers
 \end{itemize}
\end{frame}

\begin{frame}
\frametitle{Setting up the model: Income tax benchmark} 
Setup of individual problems
\begin{itemize}\setlength\itemsep{3pt}
\item We want high-wage individual $A$ pay taxes to finance transfer to beneficiary $B$
\item Each consume $C$ out of earned income $Y$ and transfer $T$, thus $C=Y+T$
\item Each exert effort $E$, defined as earned income $Y$ divided by wage $W$
\item Both individuals have same utility function $U(C,E)$
\item Government only observes income, and cannot break it down to wage and effort
\end{itemize}\medskip
Solving the problem
\begin{itemize}\setlength\itemsep{3pt}
\item If $Y > \bar{Y}$, pay tax $T$ (otherwise, receive transfer $T$)
\item But $A$ may choose to mimick $B$: Instead of $Y_A^*$, report $\bar{Y}$ (to receive transfer)
\item Hence the incentive compatibility constraint (Mimick by reducing effort, not ability)
\[
U(Y_A^*-T, Y_A^*/W_A)> U(\bar{Y}+T, \bar{Y}/W_A)
\]
\end{itemize}
\end{frame}

\begin{frame}
  \frametitle{In income tax scheme, better to restrict eligiblity}
  \begin{columns}
    \begin{column}{0.45\textwidth}
  \begin{figure}
    \centering
    \includegraphics[keepaspectratio, width=\textwidth]{nz1982_f1.png}
  \end{figure}    
    \end{column}
    \begin{column}{0.55\textwidth}
      Two utilities for $A$
      \begin{itemize}\setlength\itemsep{7pt}
        \item $\hat{U}_A$ is utility of $A$ mimicking $B$, $U_A$ is utility of being truthful
        \item If $\bar{Y}=Y_B^*$, then $A$ is indifferent between truthfully earning $Y_A^*$ and mimicking $\bar{Y}$
      \end{itemize}\medskip
      Slightly lower $\bar{Y}$ to $\bar{Y}^1$
      \begin{itemize}\setlength\itemsep{7pt}
        \item $U_B$ is reduced slightly: Productive efficiency $\downarrow$
        \item $\hat{U}_A$ significantly falls at $\bar{Y}^1$ $\to$ Better off at $Y_A^*$
        \item Room to increase $T$ without $A$ mimicking and compensate loss in $U_B$: Targeting efficiency $\uparrow$
      \end{itemize}
    \end{column}
  \end{columns}
  
\end{frame}


\begin{frame}
\frametitle{In-kind Transfers} 
Is it optimal to give in-kind transfers instead of cash equivalent?
\begin{itemize}\setlength\itemsep{3pt}
\item In general, people (and economists) prefer cash transfers over in-kind
\begin{itemize}\setlength\itemsep{3pt}
\item This is true when demand only depends on price and income
\item Exceptions: Commodity-specific externalities, paternalism etc
\end{itemize}
\item IRL, demand may also be determined by leisure, ability, medical conditions etc
\begin{itemize}\setlength\itemsep{3pt}
\item Even when income and price is held fixed, demand differs with the unobservable traits
\item Medical care consumption higher for those with poor health (even w/ income controlled)
\end{itemize}
\item We refer to these goods as \textbf{indicator goods}
\end{itemize}\medskip
What are indicator goods in reality?
\begin{itemize}\setlength\itemsep{3pt}
\item Necessary staples (rice, cooking oils)
\item Low-quality version of a good (e.g. housing)
\item Health services 
\end{itemize}
\end{frame}

\begin{frame}
  \frametitle{In-kind transfers can be used as deterrence}
  \begin{columns}
    \begin{column}{.45\textwidth}
  \begin{figure}
    \centering
    \includegraphics[keepaspectratio, width=\textwidth]{nz1982_f2.png}
  \end{figure}    
    \end{column}
    \begin{column}{.55\textwidth}
      Consumption of indicator good
      \begin{itemize}\setlength\itemsep{5pt}
        \item Out of B's income $\bar{Y}+T$, spend $X_B^*$
        \item Mimicking A spend $\hat{X}_A^*$ with the same income
      \end{itemize}\medskip
      Change the transfer scheme
      \begin{itemize}\setlength\itemsep{5pt}
        \item $\bar{X}$ is given in-kind, rest $T-\bar{X}$ as cash
        \item If $\bar{X}>\hat{X}_A^*$, $\hat{U}_A\downarrow$, and $U_B\uparrow$ until $\bar{X}=X_B^*$
        \item Better to make $\bar{X}>X_B^*$: B's marginal utility loss is small, but that of A is large
        \item This maximizes deterrence
      \end{itemize}
    \end{column}
  \end{columns}
  
\end{frame}



\begin{frame}
\frametitle{Ordeals and discussion} 
What are ordeals in this context?
\begin{itemize}\setlength\itemsep{3pt}
\item Deadweight loss that comes from eligilbity restrictions
\item Think qualification tests, administrative procedures, work requirements
\item Effective costs from ordeals vary inversely with benefits
\item This is because it works as a sorting mechanism
\end{itemize}\medskip
Discussion
\begin{itemize}\setlength\itemsep{3pt}
\item First-best solution is generally not possible (unobservable ability)
\item Here, sacrificing the benefits for intended recipient may improve social welfare
\item While it reduces welfare for the poor, it prevents unintended recipients
\item However, how much of this is acceptable still remains an open debate
\end{itemize}
\end{frame}




\subsection{Kleven, Kopczuk (2011): Transfer Program Complexity and the Take-Up of Social Benefits }
\begin{transitionframe}
\begin{center}
  \begin{huge}
    \textcolor{NTHU1}{%
Kleven, Kopczuk (2011):\\ Transfer Program Complexity and the Take-Up of Social Benefits
    }
  \end{huge}
\end{center}
\end{transitionframe}


\begin{frame}
\frametitle{Q: Optimal trade-off between complexity, eligibility, and generosity?}
Different dimensions of social programs
\begin{itemize}\setlength\itemsep{3pt}
\item Targeting (eligibility): Who is the intended recipient?
\item Complexity: What hassles are involved in receiving the benefit?
\begin{itemize}\setlength\itemsep{3pt}
\item Eligibility criteria, document requirements, trips to the office etc. 
\end{itemize}
\item Take-up: How many eligible beneficiaries do not receive the benefit?
\begin{itemize}\setlength\itemsep{3pt}
\item Significant variation across different programs (Medicare vs disability assistance)
\end{itemize}
\end{itemize}\medskip
This paper provides a framework balancing these dimensions in designing welfare program
\begin{itemize}\setlength\itemsep{3pt}
\item Consider complexity as a program feature, rather than a side effect
\item Complex programs w/ incomplete take-up can sometimes be desirable
\item Study trade-off in false rejection (Type 1) and false awards (Type 2) error framework
\begin{itemize}\setlength\itemsep{3pt}
\item Screen out undeserving recipients vs not deterring (and dropping) deserving applicants
\end{itemize}
\end{itemize}
\end{frame}

\begin{frame}
\frametitle{Setting up the model} 
Definitions of classification errors
\begin{itemize}\setlength\itemsep{3pt}
\item Type 1a (1b) error: Eligible individuals do not apply (apply but wrongly rejected)
\item Type 2 error: Ineligible individuals apply and accepted
\end{itemize}\medskip
Other underlying parameters for individuals
\begin{itemize}\setlength\itemsep{3pt}
\item Each have ability $a$, with $\sigma$ capturing precision by which this can be obsesrved
\item There is an indicator $\tilde{a}$, a signal of $a$ defined as
\[
\tilde{a} = a + \epsilon/\alpha \ \ \ \epsilon\sim N(0,\sigma^2)\implies \epsilon/\sigma \sim N(0,1)
\]
where $\epsilon$ is a noise $\alpha$ represent complexity
\item If eligilbility criteria is $\bar{a}$ and $P(\cdot)$ represent distribution of $\epsilon/\sigma$, chance of receipt is
\[
a + \epsilon/\alpha < \bar{a} \implies \epsilon / \sigma < \frac{\alpha (\bar{a}-a)}{\sigma}\implies P\left(\frac{\alpha (\bar{a}-a)}{\sigma}\right) = \bar{P}
\]
\end{itemize}
\end{frame}


\begin{frame}
\frametitle{Modeling decision to apply} 
Individual's decision to apply
\begin{itemize}\setlength\itemsep{3pt}
\item Utility depends on $a$, benefit $B$, and cost $f(\alpha)$
\item We can derive cutoff probability $\tilde{P}$ that determines application
\begin{footnotesize}
\[
\bar{P} u(a+B-f(\alpha)) + (1-\bar{P}) u (a-f(\alpha)) > u(a) \iff \bar{P} > \frac{u(a)-u(a-f(\alpha))}{u(a+B-f(\alpha))-u(a-f(\alpha))} \equiv \tilde{P}_a (\alpha, B)
\]
\end{footnotesize}
\item Individual applies when $\bar{P}>\tilde{P}$
\item Rise in $\bar{a}$ always increases application, rise in $a$ does the opposite
 \item Higher benefit $B$ lowers $\tilde{P}$ (easier to justify applying)
\item When $\bar{a}>a$, rise in $\alpha$ (complexity) and fall in $\sigma$ (better precision) increase $\bar{P}$
\item When $\bar{a}<a$ opposite conditions on $\alpha$ and $\sigma$ increase $\bar{P}$
\end{itemize}\medskip


\end{frame}


\begin{frame}
\frametitle{Government problem} 
Starting point 
\begin{itemize}\setlength\itemsep{3pt}
\item At first, those with zero type 2 and type 1a error seems ideal
\item But there is type 1b error: Imperfect testing $\to$ eligible ones may be wrongly rejected
\item We can reduce 1b error by increasing complexity, but increase 1a error
\item Raising benefits can mitigate 1a error, but induces high fiscal burden
\end{itemize}\medskip
Government goals
\begin{itemize}\setlength\itemsep{3pt}
\item Provide at least minimum income $\bar{B}$ to as many eligible population on limited budget
\begin{gather*}
  \max_{\alpha, \bar{a}, B} N_L(\alpha, \bar{a}, B) \\
  \text{s.t. } [N_L(\alpha, \bar{a}, B)+N_H(\alpha, \bar{a}, B)]B\leq R \ \ \text{and } B\geq \bar{B}
\end{gather*}
\item $L$ types have low ability (otherwise $H$ type)
\item Not every $L$ types (intended recipients) get suport, and some $H$ types get support
\end{itemize}\medskip
\end{frame}



\begin{frame}
\frametitle{Key result 1: If ability is identified precisely}
Suppose we observe both $a$ and $\sigma$
\begin{itemize}\setlength\itemsep{3pt}
\item Full separation is possible in principle: No $H$ types and only $L$ types receive benefit
\item If budget is small, provide $\bar{B}$ to $L$ groups with $B=\bar{B}$, $a_L<\bar{a}\leq a_H$
\item If very large, set $\bar{a}$ large and move towards universal provision and $B\geq \bar{B}$
\end{itemize}\medskip
If budget is intermediate (first-best not achievable):
\begin{itemize}\setlength\itemsep{3pt}
\item Full separation requires $B>\bar{B}$, $\bar{a}=a_H$, and high $\alpha$ (Proposition 2)
    \begin{itemize}\setlength\itemsep{3pt}
        \item $\alpha\uparrow$ to screen out $H$ types +  $B\uparrow$  to compensate $L$ types for the complexity (1a, 1b error $\uparrow$)
    \end{itemize}
    \item But high $\alpha$ $\to$ $L$ types deterred or wrongly rejected
    \item Forgoing full separation (let some $H$ types in) strictly 
    improves outcomes (Lemma 3)
    \begin{itemize}\setlength\itemsep{3pt}
        \item Lower $\alpha\to$ $L$ types less discouraged + $B$ no longer high to compensate for complexity
        \item Savings from lower $B$ outweigh the cost of the additional $H$ type recipients
        \item More $L$ types can be served overall (at low benefits)
    \end{itemize}
\end{itemize}
\end{frame}


\begin{frame}
\frametitle{Key result 2: Optimal programs feature all three errors}
In intermediate budget, government has three tools to screen $L$ from $H$ types: $\bar{a}$, $\alpha$, $B$
\begin{itemize}\setlength\itemsep{6pt}
    \item \textbf{Set eligibility $\bar{a}$ leniently}: $H$ types allowed, but prevents discouragement of $L$ types
    \begin{itemize}\setlength\itemsep{3pt}
        \item Since full separation is no longer the goal, $B$ can be set lower
        \item Trades off: fewer Type 1a/1b errors, but more Type 2 errors
    \end{itemize}\medskip
    \item \textbf{Use high complexity $\alpha$} to deter $H$ types from applying
    \begin{itemize}\setlength\itemsep{3pt}
        \item But this unavoidably discourages $L$ types too
        \item Trades off: fewer Type 2 errors, but more Type 1a and 1b errors
    \end{itemize}\medskip
    \item \textbf{Raise benefits $B$} to persuade deterred $L$ types back in to reduce type 1a error
    \begin{itemize}\setlength\itemsep{3pt}
        \item Only worth doing when $\bar{a}$ is set leniently - so that both $H$ \textit{and} $L$ types are induced to join
        \item Must be balanced against the budget cost of a higher $B$
    \end{itemize}
\end{itemize}\medskip

\textbf{Takeaway}: All three errors coexist in the optimum\\ ($\because$ each tool that reduces one type of error worsens another.)

\end{frame}

\begin{frame}
\frametitle{Key Takeaways}
High complexity and incomplete take-up can be \textit{optimal} policy outcomes
\begin{itemize}\setlength\itemsep{3pt}
\item Complexity is a screening tool, not merely a deadweight cost
\item All three classification errors (Type 1a, 1b, Type 2) coexist in optimal programs
    \item Eliminating one type of error worsens another: trade-offs are unavoidable
\end{itemize}\medskip
Policy implication: Is it good to keep transaction costs high?
\begin{itemize}\setlength\itemsep{3pt}
 \item Reducing transaction costs can increase take-up (we will see later)
  \item But the government may \textit{rationally} have chosen not to do so
  \item Screening benefits of complexity may outweigh the cost of incomplete take-up
  \item This justification weakens when factoring utility cost that complexity imposes
    \begin{itemize}\setlength\itemsep{3pt}
      \item Likely discourage deserving applicants (who are also poor)
      \item Note that no utility maximization is carried out here
    \end{itemize}
\end{itemize}
\end{frame}




\section{Empirical Evidence on targeting redistribution}
\subsection{Alatas, Banerjee,Hanna, Olken, Tobias (2012): Targeting the Poor: Evidence from a Field Experiment in Indonesia}
\begin{transitionframe}
\begin{center}
  \begin{huge}
    \textcolor{NTHU1}{%
Alatas, Banerjee,Hanna, Olken, Tobias (2012): \\Targeting the Poor: Evidence from a Field Experiment in Indonesia
    }
  \end{huge}
\end{center}
\end{transitionframe}



\begin{frame}
\frametitle{Q: How do different targeting mechanisms identify the beneficiaries?} 
Targeted social safety nets are increasingly becomming common
\begin{itemize}\setlength\itemsep{3pt}
\item In developed countries, this is usually done through means testing
\item Many developing countries lack credible records to implement such mechanism 
\item Targeting that does not rely on observable income is in place
\begin{itemize}\setlength\itemsep{3pt}
  \item Proxy means test (PMT): Gathers info on assets and demographics that proxy income 
  \item Community-based targeting (CBT): Beneficiaries are selected by community members
\end{itemize}
\item Prevent elite capture in communities vs better utilize local knowledge
\end{itemize}\medskip
This paper uses randomized experiments to examine which methods work better
\begin{itemize}\setlength\itemsep{3pt}
  \item Compare PBT with different extent of CBT on 640 villages in Indonesia
  \item Measure performance on per capita expenditures and satisfaction with targeting
  \item Examine where the errors in judgement is from
\end{itemize}
\end{frame}



\begin{frame}
\frametitle{Institutional context} 
Cash transfer program in Indonesia (Bantuan Langsung Tunai)
\begin{itemize}\setlength\itemsep{3pt}
\item Largest cash transfer in the developing world, started in `05 and renewed in `08
\item Targeting done based on PMT and CBT methods
\item Known to be targeted badly: 45\% of nonpoor receives it, 47\% of poor excluded
\item Many were dissatisfied with the beneficiary lists (protests)

\end{itemize}\medskip
Details of the experiment
\begin{itemize}\setlength\itemsep{3pt}
\item Provide INR 30,000 ($=\$3$) to households below location-specific poverty line
\item Government implemented the PMT in 1/3, CBT in 1/3, and hybrid in the rest
\item In CBT villages, govt surveyed only selected househods for means testing
\item Overall number of beneficiaries set in advanced, revealed later
\end{itemize}
\end{frame}

\begin{frame}
  \frametitle{Experimental design}
  Main intervention
\begin{itemize}\setlength\itemsep{3pt}
\item T1: PMT villages - based on 49 indicators (housing attributes, assets, education etc) 
\item T2: CBT villages - Community meeting to rank hh by poverty, quota revealed later
\item T3: Hybrid villages - Community ranking + verified by government officials
\item Stratified randomization based on geography
\end{itemize}\medskip
Subtreatments within CBT: Why do CBT differ from PMT?
\begin{itemize}\setlength\itemsep{3pt}
\item Elite capture: Compare mechanism with all community members vs elites only
\item Effort: Randomize order in which households are ranked to test if fatigue hurts efforts
\item Preferences: Ranks determined by broadly shared preferences or particular ones
\begin{itemize}
\item Vary meeting times (morning - women, evening - men, some meeting emphasize poverty)
\end{itemize}
\end{itemize}
\end{frame}

\begin{frame}
\frametitle{Errors included in CBT and hybrid treatments} 
\begin{figure}
  \centering
  \includegraphics[keepaspectratio, width=0.9\textwidth]{abhot2012_t1.png}
\end{figure}
\end{frame}

\begin{frame}
\frametitle{CBT does poor in capturing households just near the cutoff} 
\begin{figure}
  \centering
  \includegraphics[keepaspectratio, width=0.9\textwidth]{abhot2012_f1.png}
\end{figure}
\end{frame}


\begin{frame}
\frametitle{Elite capture does not(!) explain error rates} 
\begin{figure}
  \centering
  \includegraphics[keepaspectratio, width=0.69\textwidth]{abhot2012_t2.png}
\end{figure}
\end{frame}



\begin{frame}
\frametitle{Key takeaways and discussion} 
So what is happening with the CBT methods?
\begin{itemize}\setlength\itemsep{3pt}
\item Preference also does not play a role: no difference in outcomes due to meeting timing
  \item Communities apply standard of welfare different from per capita consumption
\item Instead, villages use different metrics
\begin{itemize}\setlength\itemsep{3pt}
  \item Education, widowhood, presence of members not in the household, household size
  \item Make guesses on households' ability to smooth consumption based on observables
\end{itemize}
\end{itemize}\medskip
Discussions
\begin{itemize}\setlength\itemsep{3pt}
  \item Little reason to advocate hybrid: Worse of the two worlds
  \item PMT matches poverty based on consumption better
  \item CBT captures measures of ability to absorb shocks not reflected in consumption
  \item Satisfication in targeting higher in CBT - result of reciprocity or other channels?
  \end{itemize}
\end{frame}

\subsection{Niehaus, Atanassova, Bertrand, Mullainathan (2013): Targeting with Agents}
\begin{transitionframe}
\begin{center}
  \begin{huge}
    \textcolor{NTHU1}{%
Niehaus, Atanassova, Bertrand, Mullainathan (2013):\\ Targeting with Agents
    }
  \end{huge}
\end{center}
\end{transitionframe}


\begin{frame}
\frametitle{Q: How to achieve optimal targeting under corruptible agency?} 
Proxy means test (PMT) is used in places where reliable data on income is absent
\begin{itemize}\setlength\itemsep{3pt}
 \item Ideally: Marginal reduction in poverty = Marginal admin cost
 \item However, many PMT departs from this rule in practice
 \item Leading explanation behind such deviation is corruption
\begin{itemize}\setlength\itemsep{3pt}
\item Diverting transfers from intended recipients
\item Inflate claims about program participation
\item Demand bribes to issue permits to eligible or ineligible recipients
\end{itemize}
\end{itemize}\medskip
This paper provides a framework on targeting rule with possible corruption
\begin{itemize}\setlength\itemsep{3pt}
  \item Principal define a targeting rule to be implemented by an agent
  \item Agents may have allocative preference, and possibly prefer bribes
  \item Principal cannot fully discipline these agents
  \item The paper provides properties of optimal targeting rules in this context
 \item Finds that more conditioning may \textit{worsen} targeting
  \end{itemize}
\end{frame}



\begin{frame}
\frametitle{Model Setup: Principal, Agent, and Households}

Households are characterized by:
\begin{itemize}\setlength\itemsep{3pt}
    \item Income $y_i \in \{\underline{y}, \bar{y}\}$ (poor or rich) and observable characteristics $\mathbf{x}_i \in X$ (e.g. own TV?)
    \item Willingness to pay $v_i$ for a program slot
\end{itemize}\medskip

Principal sets a targeting rule $R \subseteq X$:
\begin{itemize}\setlength\itemsep{3pt}
    \item Household $i$ is eligible iff $\mathbf{x}_i \in R$ ($R$: Eligible if no TV and no brick walls  )
    \item Cannot observe income directly; must rely on proxy characteristics
    \item Has progressive preferences: values transfers to poor more than rich
\end{itemize}\medskip

Agent implements the rule and may deviate:
\begin{itemize}\setlength\itemsep{3pt}
    \item Observes $(y_i, \mathbf{x}_i)$ and sets household-specific bribe prices $p_i \geq 0$
    \item Has own allocative preferences over poor and rich households
    \item Rule violations detected with probability $\pi$ and punished by fine $f$
    \item Principal cannot make $f$ arbitrarily large --- enforcement is bounded
\end{itemize}
\end{frame}




\begin{frame}
\frametitle{More Targeting Can Backfire: The Key Intuition}

Adding an eligibility indicator has \textit{two} effects that pull in opposite directions:\medskip

\begin{itemize}\setlength\itemsep{6pt}
    \item[\textbf{(+)}] \textbf{Statistical effect}: the rule better separates poor from rich households
    \item[\textbf{(-)}] \textbf{Enforcement effect}: verifying \textit{ineligibility} becomes harder when it requires checking more criteria --- officials grow less afraid of breaking the rule (i.e. become corrupt)
\end{itemize}\medskip

\textbf{Concrete example}: Suppose households with paved floors are ineligible.
\begin{itemize}\setlength\itemsep{3pt}
    \item A third party can verify ineligibility by simply observing the floor $\Rightarrow$ official is deterred
    \item Now refine the rule: ineligible only if paved floor \textit{and} a television
    \item Statistically better, but verifying ineligibility now requires checking \textit{two} things
    \item Official is less deterred $\Rightarrow$ more rule violations among ineligible households
\end{itemize}\medskip

If the enforcement effect is strong enough, it can \textit{more than offset} the statistical gain.

\end{frame}

\begin{frame}
\frametitle{When Conditions Fail: Simpler Rules Can Dominate}
Typical setup in capable countries
\begin{itemize}\setlength\itemsep{3pt}
  \item Strong enforcement: More tools available to verify eligibility
  \item Preference aligned: Principal and agent have same incentives
  \item Thus, gap in principal and agent objectives can be ignored
\end{itemize}\medskip

When none of these hold, the implications are stark:\medskip

\begin{itemize}\setlength\itemsep{6pt}
    \item Fewer, easier-to-verify criteria $>$ statistically optimal rule
    \item Sacrifices some statistical accuracy, but gains in enforceability can compensate
    \item At the extreme, universal eligiblity can provide better welfare than accurate rule
\end{itemize}\medskip

\textbf{Key policy implication}: When implementation is corruptible, more information may hurt
\begin{itemize}\setlength\itemsep{6pt}
    \item Trade-off between statistical vs. enforcement effect
\end{itemize}\medskip


\end{frame}

\begin{frame}
\frametitle{Empirical testing: India's BPL card} 
History of public programs in India
\begin{itemize}\setlength\itemsep{3pt}
 \item Originally universal public distribution system
 \item Switched to poverty targeting in 1997 citing fiscal costs
 \item Houses are now categorized into below the poverty line (BPL) or not
 \begin{itemize}\setlength\itemsep{3pt}
 \item If BPL, given cards providing access to various public services
\end{itemize}
\end{itemize}\medskip
Eligibility to BPL cards in Karnataka (research setting)
 \begin{itemize}\setlength\itemsep{3pt}
 \item Income, telephone, automobile, gas, color TV, own land, water pump etc
 \item State runs survey to check eligiblity, implemented by village officials
 \item Researchers observe statutory eligiblity + BPL receipt across 14,000 households
 \item Finds widespread rulebreaking: 48\% of households misclassified
 \begin{itemize}\setlength\itemsep{3pt}
 \item 70\% of ineligible households have BPL
 \item 13\% of eligible households do not have BPL
\end{itemize}
\end{itemize}
\end{frame}


\begin{frame}
\frametitle{Actual distribution not progressive as statutory eligibility} 
\begin{figure}
  \centering
  \includegraphics[keepaspectratio, width=0.55\textwidth]{nabm2013_f1.png}
\end{figure}
 \begin{itemize}\setlength\itemsep{3pt}
 \item Correlation with log income: -0.55 for statutory vs -0.23 for actual
\end{itemize}
\end{frame}

\begin{frame}
\frametitle{Key findings and discussion} 
Key messages
\begin{itemize}\setlength\itemsep{3pt}
    \item Using more indicators can backfire when implementation can be corruptible
    \item Rule simplicity has value: Simpler rules are harder to manipulate
    \item In such settings, trade-offs between accurately capturing poverty and enforcing rules
\end{itemize}\medskip

Discussion
\begin{itemize}\setlength\itemsep{3pt}
    \item  The case for targeting over universalism is weaker than standard analysis suggests
    \item Particularly when enforcing complex rules are difficult
    \item Corruptible or incapable agency unravels benefits of accurate targeting
\end{itemize}

\end{frame}




\subsection{Castell, Gurgand, Imbert, Tochev (2025): Take-up of Social Benefits: Experimental Evidence from France}
\begin{transitionframe}
\begin{center}
  \begin{huge}
    \textcolor{NTHU1}{%
Castell, Gurgand, Imbert, Tochev (2025):\\ Take-up of Social Benefits: Experimental Evidence from France
    }
  \end{huge}
\end{center}
\end{transitionframe}

\begin{frame}
\frametitle{Q: How to increase take-up of welfare programs?}
Many social programs suffer from low take-up
\begin{itemize}\setlength\itemsep{3pt}
\item 25\% of eligible EITC beneficiaries do not receive them
\item France: 36\% of those eligible for out-of-work minimum income do not claim it
\item Lack of information, transaction costs, welfare stigma often cited as reasons
\item Also, does alleviating barriers attract poorer or richer claimants?
\end{itemize}\medskip
This paper uses 2 nationwide experiments to study barriers to take-up of welfare programs
\begin{itemize}\setlength\itemsep{3pt}
\item Focusing on job-seekers in France
\item Experiment 1: Meet social workers to learn about eligibility and help with applications
\item Experiment 2: Job-seekers learn about eligiblity from online simulator
\item Experiment 1 increased take-up far more than experiment 2
\begin{itemize}
  \item Suggesting that transaction cost is a bigger barrier than lack of information
\end{itemize}
\end{itemize}
\end{frame}


\begin{frame}
\frametitle{Experiment design}
Experiment 1: Meeting with social workers (RDVDE)
\begin{itemize}\setlength\itemsep{3pt}
\item Sample: 60,000 job-seekers with small or no claims to unemployment benefits
\item Randomly select half of them to meeting with social services - 21\% of them joined
\item Discuss eligibility to 15 different benefits and support in applications
\item Additional variation: Mode (In-person vs phone), Content (anti-stigma, info flyer)
\end{itemize}\medskip
Experiment 2: Internet simulator on benefit eligibility
\begin{itemize}\setlength\itemsep{3pt}
\item Implemented one year before Experiment 1
\item Sample: Randomly select among 40,000 job-seekers - 18\% of treated group joined
\item Provide personalized guidance on eligibility to same range of benefits
\item Provides information, but do not lower transaction costs (no application help)
\end{itemize}
\end{frame}

\begin{frame}
\frametitle{Additional data and regression}
Additional data
\begin{itemize}\setlength\itemsep{3pt}
\item Job seekers' demographics from Pôle Emploi and CNAF
\item Benefit receipt status from RNCPS and CNAF source
\item Meeting outcomes and phone surveys conducted with French government
\end{itemize}\medskip
Regression leverages the randomized variation in treatment assignment status
\begin{itemize}\setlength\itemsep{3pt}
\item Intent-to treat regression (purely based on assignment)
\[
Y_i = \beta T_i + \delta X_i + \pi_s + e_i
\]
\item Actual effect of attendance ($M_i$): 2SLS regression
 \[
 \begin{aligned}
 M_i &= \gamma T_i + \alpha X_i + \pi_s + u_i\\
 Y_i &=\beta_{2SLS}\widehat{M}_i + \delta X_i + \pi_s + \varepsilon_i
 \end{aligned}
 \] 
\end{itemize}
\end{frame}


\begin{frame}
\frametitle{Meetings increase takeup of benefits}
\begin{figure}
  \centering
  \includegraphics[keepaspectratio, width=1\textwidth]{cgit2025_t1.png}
\end{figure}
\end{frame}

\begin{frame}
\frametitle{Simulator has little effects}
\begin{figure}
  \centering
  \includegraphics[keepaspectratio, width=1\textwidth]{cgit2025_t2.png}
\end{figure}
\end{frame}


\begin{frame}
\frametitle{Drop in claim size $\to$ Higher share of relatively richer beneficiaries}
\begin{figure}
  \centering
  \includegraphics[keepaspectratio, width=0.65\textwidth]{cgit2025_t3.png}
\end{figure}
\end{frame}


\begin{frame}
  \frametitle{Other key findings and discussion}
  To whom were the intervention the most effective?
  \begin{itemize}\setlength\itemsep{3pt}
    \item Those who were least likely to attend meetings on their own
    \begin{itemize}
      \item Estimated using marginal treatment effects (MTE) - not covered in 1st class
  \end{itemize}
  \item Effects of meeting large for components with large application costs (items to fill)
    \begin{itemize}
      \item But not too large - health benefits require going through a completely different branch
  \end{itemize}
  \item Overall, results point to transaction cost being the biggest barrier to take-up
    \end{itemize}\medskip
  Discussion on targeting
   \begin{itemize}\setlength\itemsep{3pt}
     \item The intervention did not improve targeting (claim size smaller)
     \begin{itemize}
  \item Marginal claimants induced by meetings are relatively richer (smaller claims)
\end{itemize}
\item The very poor --- most deterred by transaction costs --- remained hard to reach
     \item Targeting is a supply-side problem from the welfare provider
     \item Whereas take-up is a demand-side problem (beneficiaries)
     \end{itemize}
\end{frame}


\begin{frame}
\frametitle{Unanswered questions} 
Targeting deserving recipients and delivering services are tricky
\begin{itemize}\setlength\itemsep{3pt}
\item If information on income and preference is available, targeting income support is ideal
\item Very rarely the case: Many misclassification errors exist
\item For developing countries: Patchy information on income and monitoring agents
\begin{itemize}\setlength\itemsep{3pt}
\item Use proxy means test or community based measures: imperfect information
\item Difficult to monitor agents who mishandle allocation
\item How to improve information and personnel capacities in this setup?
\end{itemize}
\item Even if targeting is appropriate, take-up of welfare may be low
\begin{itemize}\setlength\itemsep{3pt}
\item High transaction costs for the recipient
\item Others document lack of information and welfare stigma
\item How do you increase take-up of welfare among those eligible?
\end{itemize}\medskip
\end{itemize}
\end{frame}
%%%%%%%%%%%
\end{document}
